\chapter{Background, in full}
\label{app:background}

\noindent\emph{\Cref{ch:background} states the gap this project addresses.
This is the full review behind it, organised by theme, with each citation's
scope stated precisely rather than implied.}

\section{Supernumerary robotic limbs}

The field's founding demonstrations are Parietti and Asada's bracing limbs
for aircraft fuselage assembly, later extended to whole-body
support~\cite{parietti2016support}, which anchor against the environment so
the reaction load leaves the wearer's body rather than passing through it
(\cref{fig:bracing}). In the arrangement studied here the arms are free in
space, so the chain closes through the wearer instead -- and every reaction
load is carried by a person who is not the person commanding the motion.
This difference, not the control law, is what makes wearer clearance the
binding constraint of \cref{ch:workspace}.

\begin{figure}[htbp]
  \centering
  % ORIGINAL DIAGRAM. Not a reproduction: it is drawn from the mechanical
% argument rather than copied from a published figure, so no permission is
% required and nothing is traced. It contrasts the load path of a bracing
% supernumerary limb with the load path of the arrangement studied here.
\begin{tikzpicture}[x=1mm, y=1mm, font=\scriptsize,
  body/.style   = {draw, fill=violet!10, minimum width=13mm,
                   minimum height=30mm, rounded corners=2pt},
  limb/.style   = {draw, fill=blue!10, rounded corners=2pt,
                   minimum height=6mm, text width=22mm, align=center},
  ground/.style = {draw, fill=gray!18, minimum width=52mm, minimum height=5mm}]

% ---------- left panel: bracing, chain closed through the environment
\begin{scope}
  \node[font=\bfseries, anchor=south] at (26, 44) {(a) bracing};
  \node[body]  (b1) at (10, 20) {};
  \node[anchor=north, font=\tiny, align=center] at (10, 4) {wearer};
  \node[limb]  (l1) at (34, 30) {supernumerary\\limb};
  \node[ground](g1) at (26, -6) {};
  \node[anchor=north, font=\tiny] at (26, -9) {work surface};
  \draw[flow] (b1.east) -- (l1.west);
  \draw[flow, very thick, green!45!black] (l1.south) -- (g1.north east);
  \draw[flow, very thick, green!45!black] (b1.south) -- (g1.north west);
  \node[font=\tiny\itshape, green!45!black, align=center, text width=30mm]
    at (26, 12) {the chain closes through the environment, so the reaction
                 load leaves the body};
\end{scope}

% ---------- right panel: this thesis, chain closed through the wearer
\begin{scope}[shift={(78,0)}]
  \node[font=\bfseries, anchor=south] at (26, 44) {(b) this thesis};
  \node[body]  (b2) at (10, 20) {};
  \node[anchor=north, font=\tiny, align=center] at (10, 4) {wearer\\(does not command)};
  \node[limb]  (l2) at (36, 30) {two \SI{8.2}{\kilo\gram}\\arms, free in space};
  \node[ground](g2) at (26, -6) {};
  \node[anchor=north, font=\tiny] at (26, -9) {work surface, not contacted};
  \draw[flow] (b2.east) -- (l2.west);
  \draw[flow, very thick, red!65!black] (l2.south west) to[out=250,in=20] (b2.north east);
  \node[font=\tiny\itshape, red!65!black, align=center, text width=32mm]
    at (30, 12) {the chain closes through the wearer, so every reaction load
                 is carried by a person};
\end{scope}
\end{tikzpicture}

  \caption[Two load paths for a supernumerary limb]{Bracing (the field's
  founding configuration~\cite{parietti2014bracing}) against free-space
  mounting (this platform).}
  \label{fig:bracing}
\end{figure}

Reviews since map the design space along consistent axes -- mounting
location trades reach against the moment arm imposed on the wearer,
actuation trades bandwidth against worn
mass~\cite{yang2021srlreview,prattichizzo2021augmentation,tong2021srlreview,zhang2023srlreview}
-- and a recent quantitative survey adds that the field reports embodiment
far more often than task performance, and rarely
both~\cite{zheng2026quantitative}. Worn mass is the constraint reported
least: each Gen3 masses \SI{8.2}{\kilo\gram}, the pair with harness exceeds
\SI{17}{\kilo\gram}, and simulated wearer loads suggest the biomechanical
cost is not negligible~\cite{moon2024impact}.

\section{The operator and the wearer as different people}

Fusion mounts a robotic head and two arms on a worn frame, driven by a
remote operator~\cite{saraiji2018fusion}. Two qualifications travel with
that citation: its three modes (Direct, Induced, Enforced) are levels of
\emph{physical coupling} between the robot arms and the wearer's own, not
levels of machine autonomy -- in all three the robot decides nothing -- and
the publication is a two-page demonstration with no user study and no
performance numbers. It establishes that the architecture existed; it
cannot support a claim about how well it worked.

Work on spatial guidelines for near-body interaction ran a concealed
operator, varied apparent autonomy, and measured the wearer's physiological
arousal and self-reported trust~\cite{zhou2026proxemics}: higher apparent
autonomy made both worse. Their autonomy was a hidden human and their wearer
had no task of their own, so the result does not settle what happens when
autonomy is real and the wearer is occupied -- it establishes only that the
wearer's experience is not a free variable. Disturbance rejection for the
two-person case is partly solved: Zhang et al.\ compensate human-induced
base motion and report end-effector errors of
\SI{1.37(58)}{\milli\metre}~\cite{zhang2024motioncomp}, with the human
shoulder as a floating base -- covering the case where wearer and operator
are the same person (the base motion is partly predictable to the
controller because the person commanding it also caused it), not the case
studied here.

\section{Master interfaces for dissimilar kinematics}

Guggenheim and Asada argue for miniature physical masters on the grounds
that geometric similarity supplies feedback inherently, without a separate
haptic display~\cite{guggenheim2021haptic}. The argument covers similarity
of \emph{form}; this system's master-to-slave reach ratio is roughly one to
three, which is where it departs. Joint-space correspondence preserves
posture but needs every master joint measured accurately; task-space
correspondence decouples the chains at the cost of the slave's redundant
degree of freedom (\cref{app:development}). Pinardi et al.\ review
supplementary feedback in movement augmentation and find that feedback
improving control does not reliably improve the sense of
ownership~\cite{pinardi2023feedback} -- a caution against treating either as
a proxy for the other.

\section{Shared autonomy}

Dragan and Srinivasa formalise arbitration as policy blending: a predicted
user goal and the user's actual input combined with a weight set by
prediction confidence~\cite{dragan2013policyblending}, explicit that
blending aggressively on a wrong prediction is worse than not blending.
Javdani et al.\ recast the problem as a POMDP over the user's goal,
approximated by hindsight optimisation~\cite{javdani2015hindsight}: their
user study found assistance that does not wait for high confidence lets
operators finish faster with less input. Neither line of work was developed
for a robot mounted on a second person -- there, the cost of assisting
toward the wrong goal is a manipulator moving toward someone who did not
command it, not a slower trial.

\section{Measuring embodiment, workload and trust}

The rubber hand illusion showed ownership can be induced over an object that
is plainly not one's hand~\cite{botvinick1998rubber}; Longo et al.\
decompose the resulting experience into ownership, agency and location,
which can move independently~\cite{longo2008embodiment} -- an operator may
feel authorship of a motion without feeling the arm is part of them, and a
wearer may feel neither while bearing the consequences. Workload is
conventionally the NASA Task Load Index~\cite{hart1988nasatlx}; for a
two-person system it must be administered to both, since the two will not
load the same way. Physical exertion is better tracked with a Borg
category-ratio scale~\cite{borg1982perceived}. ANSI/ISA-101.01's central
principle for operator displays~\cite{isa101} -- colour is a scarce alarm
channel, so a normal state should be largely uncoloured -- transfers
directly to a robot console.

\section{The gap, restated precisely}

Fusion established the architecture and measured nothing. Shared autonomy
raises operator performance where the robot is attached to nobody. Raising
apparent autonomy lowers a wearer's trust where the autonomy is a concealed
human and the wearer has no task. What has not been done is to measure both
ends in one system, under the same autonomy, in the same trial -- the
engineering contribution of this thesis is a system in which that question
could be asked; \cref{ch:discussionmain} sets out what still stands between
the system as built and the full session that would ask it.
